% Generated draft: requires mathematical and source review.
\documentclass[11pt]{article}
\usepackage[margin=25mm]{geometry}
\usepackage{fontspec}
\setmainfont{Latin Modern Roman}
\usepackage{amsmath,amssymb,mathtools,unicode-math}
\setmathfont{Latin Modern Math}
\usepackage{enumitem,needspace,xurl,hyperref}
\setlength{\parindent}{0pt}
\setlength{\parskip}{5pt}
\setlength{\emergencystretch}{3em}
\newcommand{\N}{\mathbb{N}}
\newcommand{\Z}{\mathbb{Z}}
\newcommand{\Q}{\mathbb{Q}}
\newcommand{\R}{\mathbb{R}}
\newcommand{\C}{\mathbb{C}}
\newcommand{\F}{\mathbb{F}}
\newcommand{\A}{\mathbb{A}}
\newcommand{\PP}{\mathbb{P}}
\newcommand{\E}{\mathbb{E}}
\newcommand{\eps}{\varepsilon}
\newcommand{\dd}{\,\mathrm{d}}
\begin{document}
\section*{UnsolvedMath ID: 700829}
\textbf{Source problem code:} catalogue entry
\subsection{Definitions and mathematical statement}
Let \(\Lambda\) be a totally ordered abelian group, written additively, whose order is invariant under translation. A \(\Lambda\)-metric on a set \(X\) is a function \(d:X\times X\to\Lambda_{\geq 0}\) satisfying symmetry, the triangle inequality, and \(d(x,y)=0\) if and only if \(x=y\). A segment joining \(x\) to \(y\) is the image of an isometric embedding of \([0,d(x,y)]_{\Lambda}\) into \(X\) taking the endpoints to \(x\) and \(y\). A \(\Lambda\)-tree is a nonempty \(\Lambda\)-metric space in which any two points are joined by a unique segment, the intersection of two segments with a common endpoint is a segment, and the union of two segments meeting only at a common endpoint is a segment.

An action of a group \(G\) on \(X\) is understood to be by isometries. It is free if \(gx=x\) implies \(g=1\) for every \(g\in G\) and \(x\in X\). It is without inversion if no element interchanges the two endpoints of a nondegenerate segment. Following Bass, as attributed in the supplied record, call \(G\) \emph{tree-free} if there exist an ordered abelian group \(\Lambda\), a \(\Lambda\)-tree \(X\), and a free action without inversion of \(G\) on \(X\). The choices of \(\Lambda\) and \(X\) may depend on \(G\).

Say that \(G\) satisfies the \emph{maximal condition for abelian subgroups} if every ascending sequence of abelian subgroups
\[
A_1\leq A_2\leq A_3\leq\cdots\leq G
\]
stabilizes: there is \(n_0\geq 1\) such that \(A_n=A_{n_0}\) for every \(n\geq n_0\). The subgroups in this condition are arbitrary abelian subgroups; no finite generation assumption is imposed on them.

\begin{enumerate}
\item[(a)] Must every finitely generated tree-free group satisfy the maximal condition for abelian subgroups?
\item[(b)] Must every finitely presented tree-free group satisfy the maximal condition for abelian subgroups?
\end{enumerate}
Here finitely presented means admitting a presentation with finitely many generators and finitely many defining relations. These are the two questions in the supplied record; its designation of the problem as open is supplied registry information, not an independently verified assessment.
\subsection{Short English statement}
For groups admitting a free action without inversion on a tree whose distances lie in some ordered abelian group, must every ascending sequence of abelian subgroups eventually become constant? Ask this first for finitely generated groups and then for finitely presented groups.
\subsection{Sources}
\begin{itemize}
\item Local UnsolvedMath registry: permanent ID \(700829\), exact source problem code \texttt{catalogue entry}, title \emph{Kourovka Notebook Problem 12.9}; supplied upstream revision \texttt{d371ee917ac9647fbb06cd6fc529deec6f3dda82}.
\item Supplied source attribution: \emph{The Kourovka Notebook (Unsolved Problems in Group Theory)}, 12th issue (1992), Problem 12.9; also identified in the record with the 21st issue edition (2026), \texttt{arXiv:1401.0300v48}. The supplied collection description names E. I. Khukhro and V. D. Mazurov as editors.
\item Proposed by A. M. Gaglione and D. Spellman in 1992. The definition of tree-free is attributed to Bass in the supplied statement; no separate Bass reference is supplied.
\item The registry labels this problem \emph{open} and records an update timestamp of \texttt{2026-10-16T14:00:00Z}. This is supplied registry metadata, not independent source verification or a literature assessment.
\end{itemize}
\end{document}
